\documentclass[10pt,a4paper]{amsart}
\usepackage{amsmath,amssymb,amsthm,graphicx}
\usepackage{verbatim}
\newenvironment{reference}{\comment}{\endcomment}
%\newenvironment{reference}{}{}
\newenvironment{slogan}{\comment}{\endcomment}
\newenvironment{history}{\comment}{\endcomment}
% For commutative diagrams we use Xy-pic
\usepackage[all]{xy}

% We use 2cell for 2-commutative diagrams.
\xyoption{2cell}
% Theorem environments.
%
\theoremstyle{plain}
\newtheorem{theorem}[subsection]{Theorem}
\newtheorem{proposition}[subsection]{Proposition}
\newtheorem{lemma}[subsection]{Lemma}

\theoremstyle{definition}
\newtheorem{definition}[subsection]{Definition}
\newtheorem{example}[subsection]{Example}
\newtheorem{exercise}[subsection]{Exercise}
\newtheorem{situation}[subsection]{Situation}

\theoremstyle{remark}
\newtheorem{remark}[subsection]{Remark}
\newtheorem{remarks}[subsection]{Remarks}

\numberwithin{equation}{subsection}

% Macros
%
\def\lim{\mathop{\mathrm{lim}}\nolimits}
\def\colim{\mathop{\mathrm{colim}}\nolimits}
\def\Spec{\mathop{\mathrm{Spec}}}
\def\Hom{\mathop{\mathrm{Hom}}\nolimits}
\def\Ext{\mathop{\mathrm{Ext}}\nolimits}
\def\SheafHom{\mathop{\mathcal{H}\!\mathit{om}}\nolimits}
\def\SheafExt{\mathop{\mathcal{E}\!\mathit{xt}}\nolimits}
\def\Sch{\mathit{Sch}}
\def\Mor{\mathop{\mathrm{Mor}}\nolimits}
\def\Ob{\mathop{\mathrm{Ob}}\nolimits}
\def\Sh{\mathop{\mathit{Sh}}\nolimits}
\def\NL{\mathop{N\!L}\nolimits}
\def\CH{\mathop{\mathrm{CH}}\nolimits}
\def\proetale{{pro\text{-}\acute{e}tale}}
\def\etale{{\acute{e}tale}}
\def\QCoh{\mathit{QCoh}}
\def\Ker{\mathop{\mathrm{Ker}}}
\def\Im{\mathop{\mathrm{Im}}}
\def\Coker{\mathop{\mathrm{Coker}}}
\def\Coim{\mathop{\mathrm{Coim}}}

% Boxtimes
%
\DeclareMathSymbol{\boxtimes}{\mathbin}{AMSa}{"02}

%
% Macros for moduli stacks/spaces
%
\def\QCohstack{\mathcal{QC}\!\mathit{oh}}
\def\Cohstack{\mathcal{C}\!\mathit{oh}}
\def\Spacesstack{\mathcal{S}\!\mathit{paces}}
\def\Quotfunctor{\mathrm{Quot}}
\def\Hilbfunctor{\mathrm{Hilb}}
\def\Curvesstack{\mathcal{C}\!\mathit{urves}}
\def\Polarizedstack{\mathcal{P}\!\mathit{olarized}}
\def\Complexesstack{\mathcal{C}\!\mathit{omplexes}}
% \Pic is the operator that assigns to X its picard group, usage \Pic(X)
% \Picardstack_{X/B} denotes the Picard stack of X over B
% \Picardfunctor_{X/B} denotes the Picard functor of X over B
\def\Pic{\mathop{\mathrm{Pic}}\nolimits}
\def\Picardstack{\mathcal{P}\!\mathit{ic}}
\def\Picardfunctor{\mathrm{Pic}}
\def\Deformationcategory{\mathcal{D}\!\mathit{ef}}
\usepackage{cleveref}
\crefname{originallemma}{original lemma}{original lemmas}
\Crefname{originallemma}{Original lemma}{Original lemmas}
\makeatletter
\newlabel{decent-spaces-lemma-bounded-fibres}{{03JX}{source}{}{sourcepointer-1}{}}
\newlabel{decent-spaces-lemma-bounded-fibres@cref}{{[originallemma][0][]03JX}{[1][0][]source}{}{}{}}
\makeatother
\crefname{originallemma}{original lemma}{original lemmas}
\Crefname{originallemma}{Original lemma}{Original lemmas}
\makeatletter
\newlabel{decent-spaces-lemma-quasi-compact-reasonable-stratification}{{07S8}{source}{}{sourcepointer-2}{}}
\newlabel{decent-spaces-lemma-quasi-compact-reasonable-stratification@cref}{{[originallemma][0][]07S8}{[1][0][]source}{}{}{}}
\makeatother
\crefname{originallemma}{original lemma}{original lemmas}
\Crefname{originallemma}{Original lemma}{Original lemmas}
\makeatletter
\newlabel{morphisms-lemma-H-quasi-projective-open-H-projective}{{01WA}{source}{}{sourcepointer-3}{}}
\newlabel{morphisms-lemma-H-quasi-projective-open-H-projective@cref}{{[originallemma][0][]01WA}{[1][0][]source}{}{}{}}
\makeatother
\crefname{originallemma}{original lemma}{original lemmas}
\Crefname{originallemma}{Original lemma}{Original lemmas}
\makeatletter
\newlabel{morphisms-lemma-locally-projective-proper}{{01WC}{source}{}{sourcepointer-4}{}}
\newlabel{morphisms-lemma-locally-projective-proper@cref}{{[originallemma][0][]01WC}{[1][0][]source}{}{}{}}
\makeatother
\crefname{originallemma}{original lemma}{original lemmas}
\Crefname{originallemma}{Original lemma}{Original lemmas}
\makeatletter
\newlabel{morphisms-lemma-quasi-affine-finite-type-over-S}{{04II}{source}{}{sourcepointer-5}{}}
\newlabel{morphisms-lemma-quasi-affine-finite-type-over-S@cref}{{[originallemma][0][]04II}{[1][0][]source}{}{}{}}
\makeatother
\crefname{originallemma}{original lemma}{original lemmas}
\Crefname{originallemma}{Original lemma}{Original lemmas}
\makeatletter
\newlabel{morphisms-lemma-quasi-compact-immersion}{{01RG}{source}{}{sourcepointer-6}{}}
\newlabel{morphisms-lemma-quasi-compact-immersion@cref}{{[originallemma][0][]01RG}{[1][0][]source}{}{}{}}
\makeatother
\crefname{originalsection}{original section}{original sections}
\Crefname{originalsection}{Original section}{Original sections}
\makeatletter
\newlabel{morphisms-section-projective}{{01W7}{source}{}{sourcepointer-7}{}}
\newlabel{morphisms-section-projective@cref}{{[originalsection][0][]01W7}{[1][0][]source}{}{}{}}
\makeatother
\crefname{originalproposition}{original proposition}{original propositions}
\Crefname{originalproposition}{Original proposition}{Original propositions}
\makeatletter
\newlabel{schemes-proposition-characterize-universally-closed}{{01KF}{source}{}{sourcepointer-8}{}}
\newlabel{schemes-proposition-characterize-universally-closed@cref}{{[originalproposition][0][]01KF}{[1][0][]source}{}{}{}}
\makeatother
\crefname{originallemma}{original lemma}{original lemmas}
\Crefname{originallemma}{Original lemma}{Original lemmas}
\makeatletter
\newlabel{spaces-morphisms-lemma-diagonal-unramified-morphism}{{05W1}{source}{}{sourcepointer-9}{}}
\newlabel{spaces-morphisms-lemma-diagonal-unramified-morphism@cref}{{[originallemma][0][]05W1}{[1][0][]source}{}{}{}}
\makeatother
\crefname{originallemma}{original lemma}{original lemmas}
\Crefname{originallemma}{Original lemma}{Original lemmas}
\makeatletter
\newlabel{spaces-morphisms-lemma-equality-of-morphisms}{{084N}{source}{}{sourcepointer-10}{}}
\newlabel{spaces-morphisms-lemma-equality-of-morphisms@cref}{{[originallemma][0][]084N}{[1][0][]source}{}{}{}}
\makeatother
\crefname{originallemma}{original lemma}{original lemmas}
\Crefname{originallemma}{Original lemma}{Original lemmas}
\makeatletter
\newlabel{spaces-morphisms-lemma-etale-unramified}{{06CR}{source}{}{sourcepointer-11}{}}
\newlabel{spaces-morphisms-lemma-etale-unramified@cref}{{[originallemma][0][]06CR}{[1][0][]source}{}{}{}}
\makeatother
\crefname{originallemma}{original lemma}{original lemmas}
\Crefname{originallemma}{Original lemma}{Original lemmas}
\makeatletter
\newlabel{spaces-morphisms-lemma-lift-valuation-ring-through-flat-morphism}{{089F}{source}{}{sourcepointer-12}{}}
\newlabel{spaces-morphisms-lemma-lift-valuation-ring-through-flat-morphism@cref}{{[originallemma][0][]089F}{[1][0][]source}{}{}{}}
\makeatother
\crefname{originallemma}{original lemma}{original lemmas}
\Crefname{originallemma}{Original lemma}{Original lemmas}
\makeatletter
\newlabel{spaces-morphisms-lemma-refined-valuative-criterion-universally-closed}{{089G}{source}{}{sourcepointer-13}{}}
\newlabel{spaces-morphisms-lemma-refined-valuative-criterion-universally-closed@cref}{{[originallemma][0][]089G}{[1][0][]source}{}{}{}}
\makeatother
\title{A separated finite type algebraic space over an affine base admits a proper surjective cover by an H-quasi-projective scheme (the weak Chow lemma)}
\author{The Stacks Project Authors}
\date{stacks-project:spaces-cohomology.tex:local-snapshot-20261012; GNU FDL 1.2 or later}
\begin{document}
\maketitle
\noindent\textit{Editorial note.} All delivered mathematical text is the author's own, copied from the pinned Stacks Project chapter file spaces-cohomology.tex; no statement or proof step was rewritten, reordered or removed.\par\medskip
\noindent\textit{Editorial note.} Preamble provenance: the pinned chapter file opens with \textbackslash{}input\{preamble\} and carries no preamble text of its own, so the project's single Stacks-wide preamble staged at sources/stacks-spaces-cohomology/preamble.tex is excerpted verbatim as three spans, lines 8-12 (verbatim package and the project's own reference/slogan/history display environments), lines 14-18 (Xy-pic setup) and lines 154-217 (theorem declarations and mathematical macros). The wrapper class and the amsmath/amssymb/amsthm/graphicx base packages are supplied by the assembler. No macro was invented or redefined.\par\medskip
\noindent\textit{Editorial note.} Same-chapter pointers are located at their \textbackslash{}label lines inside the pinned chapter file; cross-chapter pointers are located in the project's own tag list staged in the snapshot (sources/stacks-chow/tags.txt), at the verbatim lines carrying their canonical tags.\par\medskip
\noindent\textit{Allowed standard prerequisites (editorial).}\begin{itemize}
\item Boundedness of the geometric fibres of a surjective etale cover, and the stratification where the cover becomes finite of the minimal degree. Decent Spaces, labels decent-spaces-lemma-bounded-fibres (invoked at proof line 3510) and decent-spaces-lemma-quasi-compact-reasonable-stratification (invoked at proof line 3514). They produce an integer \$d\$ bounding the number of geometric fibre points of the surjective etale cover \$W \textbackslash{}to X\$, and, for \$d\$ minimal, a nonempty open \$U \textbackslash{}subset X\$ over which the cover is finite etale of degree \$d\$.
\item The diagonal of a separated etale morphism is an open and closed immersion. Morphisms of Spaces, labels spaces-morphisms-lemma-etale-unramified and spaces-morphisms-lemma-diagonal-unramified-morphism (invoked at proof lines 3524 and 3525). It makes the diagonals open and closed subschemes of the \$d\$-fold fibre product, hence the complement \$V\$ of all diagonals is a quasi-compact scheme.
\item Every affine scheme of finite type embeds as an open subscheme of an H-projective scheme. Morphisms, labels morphisms-lemma-quasi-affine-finite-type-over-S and morphisms-lemma-H-quasi-projective-open-H-projective (invoked at proof lines 3533 and 3534). It provides the open immersion \$W \textbackslash{}subset Y\$ with \$Y\$ H-projective over \$A\$ into which the construction is placed.
\item The refined valuative criterion for universal closedness, and lifting points through the flat cover after extending the valuation ring. Morphisms of Spaces, labels spaces-morphisms-lemma-refined-valuative-criterion-universally-closed (invoked at proof line 3576) and spaces-morphisms-lemma-lift-valuation-ring-through-flat-morphism (invoked at proof line 3596). The criterion reduces properness of \$X' \textbackslash{}to X\$ to lifting valuation ring diagrams, and the lifting lemma supplies \$w : \textbackslash{}Spec(R) \textbackslash{}to W\$ over \$\textbackslash{}Spec(R) \textbackslash{}to X\$ after extending \$R\$.
\item The valuative criterion of properness for the projective morphisms \$g\_i\$. Morphisms, label morphisms-lemma-locally-projective-proper and Schemes, label schemes-proposition-characterize-universally-closed (invoked at proof lines 3606-3608). They lift \$w\$ to a morphism \$z : \textbackslash{}Spec(R) \textbackslash{}to Z\$ whose image lands in \$g\_i\textasciicircum{}\{-1\}(W) \textbackslash{}subset X'\$, completing the verification of the claim.
\end{itemize}\medskip
\ifdefined\hypertarget\hypertarget{sourcepointer-1}{}\fi
\noindent\textit{Source reference.} Original lemma 03JX: sources/stacks-chow/tags.txt, source lines 4655--4655. Navigation pointer for \textbackslash{}ref\{decent-spaces-lemma-bounded-fibres\} used inside the counted span (line 3510); the argument that consumes it is retained in the span.\par\medskip
\ifdefined\hypertarget\hypertarget{sourcepointer-2}{}\fi
\noindent\textit{Source reference.} Original lemma 07S8: sources/stacks-chow/tags.txt, source lines 9893--9893. Navigation pointer for \textbackslash{}ref\{decent-spaces-lemma-quasi-compact-reasonable-stratification\} used inside the counted span (line 3514); the argument that consumes it is retained in the span.\par\medskip
\ifdefined\hypertarget\hypertarget{sourcepointer-3}{}\fi
\noindent\textit{Source reference.} Original lemma 01WA: sources/stacks-chow/tags.txt, source lines 2475--2475. Navigation pointer for \textbackslash{}ref\{morphisms-lemma-H-quasi-projective-open-H-projective\} used inside the counted span (line 3534); the argument that consumes it is retained in the span.\par\medskip
\ifdefined\hypertarget\hypertarget{sourcepointer-4}{}\fi
\noindent\textit{Source reference.} Original lemma 01WC: sources/stacks-chow/tags.txt, source lines 2477--2477. Navigation pointer for \textbackslash{}ref\{morphisms-lemma-locally-projective-proper\} used inside the counted span (line 3609); the argument that consumes it is retained in the span.\par\medskip
\ifdefined\hypertarget\hypertarget{sourcepointer-5}{}\fi
\noindent\textit{Source reference.} Original lemma 04II: sources/stacks-chow/tags.txt, source lines 5904--5904. Navigation pointer for \textbackslash{}ref\{morphisms-lemma-quasi-affine-finite-type-over-S\} used inside the counted span (line 3533); the argument that consumes it is retained in the span.\par\medskip
\ifdefined\hypertarget\hypertarget{sourcepointer-6}{}\fi
\noindent\textit{Source reference.} Original lemma 01RG: sources/stacks-chow/tags.txt, source lines 2301--2301. Navigation pointer for \textbackslash{}ref\{morphisms-lemma-quasi-compact-immersion\} used inside the counted span (line 3545); the argument that consumes it is retained in the span.\par\medskip
\ifdefined\hypertarget\hypertarget{sourcepointer-7}{}\fi
\noindent\textit{Source reference.} Original section 01W7: sources/stacks-chow/tags.txt, source lines 2472--2472. Navigation pointer for \textbackslash{}ref\{morphisms-section-projective\} used inside the counted span (line 3548); the argument that consumes it is retained in the span.\par\medskip
\ifdefined\hypertarget\hypertarget{sourcepointer-8}{}\fi
\noindent\textit{Source reference.} Original proposition 01KF: sources/stacks-chow/tags.txt, source lines 2048--2048. Navigation pointer for \textbackslash{}ref\{schemes-proposition-characterize-universally-closed\} used inside the counted span (line 3611); the argument that consumes it is retained in the span.\par\medskip
\ifdefined\hypertarget\hypertarget{sourcepointer-9}{}\fi
\noindent\textit{Source reference.} Original lemma 05W1: sources/stacks-chow/tags.txt, source lines 7570--7570. Navigation pointer for \textbackslash{}ref\{spaces-morphisms-lemma-diagonal-unramified-morphism\} used inside the counted span (line 3525); the argument that consumes it is retained in the span.\par\medskip
\ifdefined\hypertarget\hypertarget{sourcepointer-10}{}\fi
\noindent\textit{Source reference.} Original lemma 084N: sources/stacks-chow/tags.txt, source lines 10328--10328. Navigation pointer for \textbackslash{}ref\{spaces-morphisms-lemma-equality-of-morphisms\} used inside the counted span (line 3557); the argument that consumes it is retained in the span.\par\medskip
\ifdefined\hypertarget\hypertarget{sourcepointer-11}{}\fi
\noindent\textit{Source reference.} Original lemma 06CR: sources/stacks-chow/tags.txt, source lines 8155--8155. Navigation pointer for \textbackslash{}ref\{spaces-morphisms-lemma-etale-unramified\} used inside the counted span (line 3524); the argument that consumes it is retained in the span.\par\medskip
\ifdefined\hypertarget\hypertarget{sourcepointer-12}{}\fi
\noindent\textit{Source reference.} Original lemma 089F: sources/stacks-chow/tags.txt, source lines 10495--10495. Navigation pointer for \textbackslash{}ref\{spaces-morphisms-lemma-lift-valuation-ring-through-flat-morphism\} used inside the counted span (line 3597); the argument that consumes it is retained in the span.\par\medskip
\ifdefined\hypertarget\hypertarget{sourcepointer-13}{}\fi
\noindent\textit{Source reference.} Original lemma 089G: sources/stacks-chow/tags.txt, source lines 10496--10496. Navigation pointer for \textbackslash{}ref\{spaces-morphisms-lemma-refined-valuative-criterion-universally-closed\} used inside the counted span (line 3577); the argument that consumes it is retained in the span.\par\medskip
\begin{lemma}
\label{lemma-weak-chow}
Let $A$ be a ring. Let $X$ be an algebraic space over $\Spec(A)$
whose structure morphism $X \to \Spec(A)$ is separated of finite type.
Then there exists a proper surjective morphism $X' \to X$
where $X'$ is a scheme which is H-quasi-projective over $\Spec(A)$.
\end{lemma}

\begin{proof}
Let $W$ be an affine scheme and let $f : W \to X$ be a surjective
\'etale morphism. There exists an integer $d$ such that all geometric
fibres of f have $\leq d$ points (because $X$ is a separated algebraic
hence reasonable, see
Decent Spaces, Lemma \ref{decent-spaces-lemma-bounded-fibres}).
Picking $d$ minimal we get a nonempty open $U \subset X$ such that
$f^{-1}(U) \to U$ is finite \'etale of degree $d$, see
Decent Spaces, Lemma
\ref{decent-spaces-lemma-quasi-compact-reasonable-stratification}.
Let
$$
V \subset W \times_X W \times_X \ldots \times_X W
$$
($d$ factors in the fibre product) be the complement of all the diagonals.
Because $W \to X$ is separated the diagonal $W \to W \times_X W$ is a
closed immersion. Since $W \to X$ is \'etale the diagonal
$W \to W \times_X W$ is an open immersion, see
Morphisms of Spaces, Lemmas
\ref{spaces-morphisms-lemma-etale-unramified} and
\ref{spaces-morphisms-lemma-diagonal-unramified-morphism}.
Hence the diagonals are open and closed subschemes
of the quasi-compact scheme $W \times_X \ldots \times_X W$.
In particular we conclude $V$ is a quasi-compact scheme.
Choose an open immersion $W \subset Y$ with $Y$ H-projective over
$A$ (this is possible as $W$ is affine and of finite type over $A$;
for example we can use
Morphisms, Lemmas
\ref{morphisms-lemma-quasi-affine-finite-type-over-S} and
\ref{morphisms-lemma-H-quasi-projective-open-H-projective}).
Let
$$
Z \subset Y \times_A Y \times_A \ldots \times_A Y
$$
be the scheme theoretic image of the composition
$V \to W \times_X \ldots \times_X W \to Y \times_A \ldots \times_A Y$.
Observe that this morphism is quasi-compact since $V$ is quasi-compact
and $Y \times_A \ldots \times_A Y$ is separated.
Note that $V \to Z$ is an open immersion as
$V \to Y \times_A \ldots \times_A Y$ is an immersion, see
Morphisms, Lemma \ref{morphisms-lemma-quasi-compact-immersion}.
The projection morphisms give $d$ morphisms $g_i : Z \to Y$.
These morphisms $g_i$ are projective as $Y$ is projective over $A$, see
material in Morphisms, Section \ref{morphisms-section-projective}.
We set
$$
X' = \bigcup g_i^{-1}(W) \subset Z
$$
There is a morphism $X' \to X$ whose restriction to $g_i^{-1}(W)$ is
the composition $g_i^{-1}(W) \to W \to X$.
Namely, these morphisms agree over $V$ hence agree over
$g_i^{-1}(W) \cap g_j^{-1}(W)$ by
Morphisms of Spaces, Lemma \ref{spaces-morphisms-lemma-equality-of-morphisms}.
Claim: the morphism $X' \to X$ is proper.

\medskip\noindent
If the claim holds, then the lemma follows by induction on $d$.
Namely, by construction $X'$ is H-quasi-projective over $\Spec(A)$.
The image of $X' \to X$ contains the open $U$ as $V$ surjects onto $U$.
Denote $T$ the reduced induced algebraic space structure on $X \setminus U$.
Then $T \times_X W$ is a closed subscheme of $W$, hence affine.
Moreover, the morphism $T \times_X W \to T$ is \'etale and every geometric
fibre has $< d$ points. By induction hypothesis there exists a proper
surjective morphism $T' \to T$ where $T'$ is a scheme H-quasi-projective
over $\Spec(A)$. Since $T$ is a closed subspace of $X$ we see that
$T' \to X$ is a proper morphism. Thus the lemma follows by taking the
proper surjective morphism $X' \amalg T' \to X$.

\medskip\noindent
Proof of the claim. By construction the morphism $X' \to X$ is separated
and of finite type. We will check conditions (1) -- (4) of
Morphisms of Spaces, Lemma
\ref{spaces-morphisms-lemma-refined-valuative-criterion-universally-closed}
for the morphisms $V \to X'$ and $X' \to X$.
Conditions (1) and (2) we have seen above.
Condition (3) holds as $X' \to X$ is separated (as a morphism whose
source is a separated algebraic space). Thus it suffices to check
liftability to $X'$ for diagrams
$$
\xymatrix{
\Spec(K) \ar[r] \ar[d] & V \ar[d] \\
\Spec(R) \ar[r] & X
}
$$
where $R$ is a valuation ring with fraction field $K$.
Note that the top horizontal map is given by $d$ pairwise distinct
$K$-valued points $w_1, \ldots, w_d$ of $W$. In fact, this
is a complete set of inverse images of the point $x \in X(K)$
coming from the diagram. Since $W \to X$ is surjective,
we can, after possibly replacing $R$ by an extension of valuation rings,
lift the morphism $\Spec(R) \to X$ to a morphism $w : \Spec(R) \to W$, see
Morphisms of Spaces, Lemma
\ref{spaces-morphisms-lemma-lift-valuation-ring-through-flat-morphism}.
Since $w_1, \ldots, w_d$ is a complete collection of inverse images of
$x$ we see that $w|_{\Spec(K)}$ is equal to one of them, say $w_i$.
Thus we see that we get a commutative diagram
$$
\xymatrix{
\Spec(K) \ar[r] \ar[d] &  Z \ar[d]_{g_i}\\
\Spec(R) \ar[r]^w & Y
}
$$
By the valuative criterion of properness for the projective
morphism $g_i$ we can lift $w$ to $z : \Spec(R) \to Z$, see
Morphisms, Lemma \ref{morphisms-lemma-locally-projective-proper}
and
Schemes, Proposition \ref{schemes-proposition-characterize-universally-closed}.
The image of $z$ is in $g_i^{-1}(W) \subset X'$ and the proof is complete.
\end{proof}

\end{document}
